\scanpage{44}
des ``espaces stratifiés'' (en un sens qui
va être précisé plus bas) dont les
``\unsure{pierres|pièces}'' sont des $K(\unsure{\pi|\Pi},1)$, \deleted{\unsure{cela devrait}}\inserted{tels}
les espaces modulaires de Teichmüller
et leurs compactifiés, et \unsure{les|leurs} sous-espaces
\inserted{et voisinages tubulaires parmi ceux-ci}
qu'on peut leur associer grâce à leurs
stratifications canoniques, le type
d'homotopie doit pouvoir s'exprimer
en termes d'un système inductif
convenable (le plus souvent à ens. d'in-
dices fini) de groupoïdes fondamentaux
\unsure{desdites|des dites} \unsure{pierres|pièces}.

Il me semble d'ailleurs après-coup,
réflexion faite, que cette théorie ``d'intégration''
des types d'homotopie, que nébuleusement
je voyais comme liée au grand yoga des
relations entre structures homotopiques (relatives
sur un topos disons) et les $n$-groupoïdes,
est en fait relativement triviale, et
justiciable sans doute de techniques $\frac12$-
simpliciales des plus standart. Il me
faudra revenir là dessus par la suite -
% Pale bottom-left foliation in another hand; no authorial page number at the top.
\margin{43}
